\section{The hyperplane loss identity}
\label{sec:hyperplane-loss}

The identity is a count of secant--point incidences inside a hyperplane,
with the cap points included at a fictitious multiplicity so that every
secant contributes the same way.

\begin{proof}[Proof of Theorem~\ref{thm:hyperplane-loss}]
Extend the index to the cap points by \(r(a)=k-1\) for \(a\in A\), the
number of secants through \(a\).  Every secant meets \(\Pi\) in exactly
one point unless it lies in \(\Pi\), in which case it contributes
\(q+1\) incidences, and the secants inside \(\Pi\) are the
\(\binom{s}{2}\) pairs of \(A\cap\Pi\).  Hence
\[
 \sum_{x\in\Pi}r(x)=N+q\binomtwo{s}.
\]
Subtracting one for each of the \(\theta_{d-1}\) points of \(\Pi\), and
then \(k-2\) for each of the \(s\) cap points, leaves
\(\sum_{x\in\Pi\setminus A}(r(x)-1)=N+q\binom{s}{2}-\theta_{d-1}-(k-2)s\),
which is~\eqref{eq:hyperplane-loss}.  If \(A\) is complete then every
summand \(r(x)-1\) is nonnegative, so \(Q(s)\ge0\), and the summands over
all of \(\PG(d,q)\setminus A\) add to \(\Lambda_0\)
by~\eqref{eq:overlap-loss}, so \(Q(s)\le\Lambda_0\).
\end{proof}

Summing~\eqref{eq:hyperplane-loss} over all hyperplanes and using the
character equations~\eqref{eq:character} returns
\(\theta_{d-1}\sum_{x\notin A}(r(x)-1)=\theta_{d-1}\Lambda(A)\) on both
sides, as it must; the content of the theorem is the distribution, not
the total.  For \(d=2\) the hyperplanes are lines, \(s\le2\), and the
identity reads \(Q(0)=N-q-1\), \(Q(1)=N-k-q+1\), \(Q(2)=N-2k+q+3\):
a complete arc in \(\PG(2,q)\) has \(N\ge q+1\) with the loss on an
external line equal to \(N-q-1\).  The planar case carries no new
information.  For \(d\ge3\) the quadratic \(Q\) has an interior minimum
near \(s=(k-2)/q+\tfrac12\), and the shape of the admissible set is
governed by two numbers.

\begin{lemma}[Shape of the admissible set]
\label{lem:admissible-shape}
Let \(A\) be a complete \(k\)-cap in \(\PG(d,q)\), \(d\ge3\).  If
\(Q\) has real roots \(\sigma_1\le\sigma_2\), then every hyperplane
section size lies in
\[
 \cS=\{s\le\sigma_1:Q(s)\le\Lambda_0\}\cup\{s\ge\sigma_2:Q(s)\le\Lambda_0\},
\]
two runs of consecutive integers, each with at most
\(2+\sqrt{2\Lambda_0/q}\) elements; no section size lies strictly between
the roots.  If \(Q\) has no real root, then \(\cS\) is the single run
\(\{s:Q(s)\le\Lambda_0\}\).
\end{lemma}

\begin{proof}
The vertex of \(Q\) is at \(s^*=\tfrac12+(k-2)/q\), and on each side of
it \(Q\) is monotone, so each run is an interval of integers.  For
\(s\ge s^*\) and \(j\ge0\),
\(Q(s+j)-Q(s)=j\bigl(qs-(k-2)\bigr)+q\binom{j}{2}\ge q\binom{j}{2}\),
and for \(s\le s^*\),
\(Q(s-j)-Q(s)=j\bigl((k-2)-qs\bigr)+q\binom{j+1}{2}\ge q\binom{j}{2}\).
If both endpoints of a run are admissible, the difference of their
\(Q\)-values is at most \(\Lambda_0\), so \(q\binom{j}{2}\le\Lambda_0\)
for the distance \(j\) between them.
\end{proof}

At the counting bound \(\Lambda_0<k(q-1)+1\), so each run has
\(O(\sqrt k)\) elements, while the two runs are separated by the distance
between the roots of \(Q\).  For the five parameter sets of
Theorem~\ref{thm:exclusions} the runs have at most three elements and the
gap between them is at least four; Table~\ref{tab:cases} lists them.

\subsection{Prescribed holes}

The identity is exact for every cap, so it can be applied when only part
of the space must be covered.  Let \(\cH\) be a set of \(h\) points off
\(A\), and call \(A\) complete outside \(\cH\) if every point outside
\(A\cup\cH\) is covered; the notion is studied
in~\cite{RuddSecantDefects}, where \(\cH\) is a conic, a quadric, or an
arbitrary set.

\begin{proposition}[Loss identity with holes]
\label{prop:holes}
Let \(A\) be a \(k\)-cap complete outside \(\cH\), \(|\cH|=h\), and let
\(u\le h\) be the number of uncovered points.  Then for every hyperplane
\(\Pi\),
\[
 -|\cH\cap\Pi|\le -\#\{x\in\cH\cap\Pi:r(x)=0\}\le Q(s_\Pi)\le\Lambda_0+u\le\Lambda_0+h .
\]
In particular every section size lies in
\(\{s:-h\le Q(s)\le\Lambda_0+h\}\).
\end{proposition}

\begin{proof}
In~\eqref{eq:hyperplane-loss} the summands \(r(x)-1\) are at least
\(-1\), with equality exactly at uncovered points, all of which lie in
\(\cH\); this gives the lower bounds.  The nonnegative summands over the
whole space add to \(\Lambda(A)+u\) with \(\Lambda(A)\) as
in~\eqref{eq:overlap-loss}, and \(\Lambda(A)=N(q-1)-(\theta_d-k-u)=\Lambda_0+u\).
\end{proof}

Every statement below that uses only the admissible set and the moment
identities holds for caps complete outside \(\cH\) after replacing
\(\Lambda_0\) by \(\Lambda_0+h\) and \(\cS\) by the wider set of the
proposition.  We do not pursue the numerical consequences here.
